Formal mathematics libraries are rapidly expanding, driven by human contributors and advances in neural theorem proving. While proof assistants such as Lean~4 guarantee correctness, they offer no guarantees about proof quality: proofs may be verbose, opaque, or heavily dependent on obscure library lemmas. Human maintainers cannot keep pace with the growing volume of machine-generated proofs. This thesis addresses the problem of \textit{automated proof optimization} --- rewriting a verified proof to be better with respect to a user-defined criterion while preserving formal correctness.

We present two systems. \textbf{ImProver} is the first proof optimization agent: a large language model system that rewrites Lean proofs using Chain-of-States prompting (which annotates intermediate proof states into the context), iterative error correction, and retrieval-augmented generation. ImProver outperforms naive GPT-4o by $423$--$778\%$ on optimization metrics across undergraduate, competition, and research-level mathematics, and demonstrates that proof optimization generalizes neural theorem proving.

However, ImProver relies on expensive proprietary models, limiting its scalability. We address this with \textbf{ImProver~2}, a self-improving pipeline that trains a 7B-parameter small language model via iterative preference optimization. ImProver~2 augments the model with neurosymbolic context --- goal-state traces, relevant lemma signatures, and auto-informalized proof sketches --- and uses a novel replay buffer to stabilize training across multiple rounds. The resulting model outperforms models orders-of-magnitude larger on the structural metrics of \textit{modularity} (decomposing proofs into reusable subproofs) and \textit{dependencies} (reducing explicitly named external lemmas), leading all evaluated systems on modularity and matching frontier models on dependencies.

A central finding of this thesis is that \textit{formal structure exposure} is the key ingredient for proof optimization at both inference time (CoS prompting in ImProver) and training time (neurosymbolic augmentation in ImProver~2). Together, these systems establish proof optimization as a scalable, learnable task and set the stage for further work on AI-assisted formal mathematics.
